% !TEX program = lualatex
%!TEX TS-program = lualatex
\documentclass[12pt,a4paper]{article}
\usepackage{fontspec}
\usepackage{polyglossia}
\setdefaultlanguage{polish}
\usepackage{geometry}
\geometry{margin=2.1cm}
\usepackage[table]{xcolor}
\usepackage{longtable}
\usepackage{array}
\usepackage{booktabs}
\usepackage{enumitem}
\usepackage{ragged2e}
\definecolor{tableHeader}{HTML}{E8EEF7}
\definecolor{tableStripe}{HTML}{F7F9FC}
\newcolumntype{L}[1]{>{\RaggedRight\arraybackslash}p{#1}}
\newcolumntype{C}[1]{>{\Centering\arraybackslash}p{#1}}
\renewcommand{\arraystretch}{1.22}
\setlength{\tabcolsep}{4pt}
\setlist[itemize]{leftmargin=*, itemsep=2pt, topsep=2pt}
\setlength{\parindent}{0pt}
\setlength{\parskip}{6pt}
\emergencystretch=3em
\sloppy
\begin{document}

\section*{Raport ze sprintu nr 7 - Calisia Banking App}
Data raportu: 06.06.2026

Okres sprintu: 01.06 - 06.06.2026

Product Owner / Scrum Master: Nikola Ruczyńska

\subsection*{Lista osób}
\begin{itemize}
\item Kubryn Jeremiasz
\item Łaszuk Jakub
\item Makarevich Ksenia
\item Nuckowski Michał
\item Ruczyńska Nikola
\item Wołczko Jakub
\item Olszewski Dawid
\end{itemize}

\subsection*{Product Owner i Scrum Master sprintu}
Nikola Ruczyńska

\subsection*{Lista historyjek na sprint}
\begingroup
\footnotesize
\rowcolors{2}{tableStripe}{white}
\begin{longtable}{@{}C{0.08\textwidth}L{0.62\textwidth}C{0.10\textwidth}L{0.12\textwidth}@{}}
\toprule
\rowcolor{tableHeader}
\textbf{Kod} & \textbf{Historyjka} & \textbf{SP} & \textbf{Status} \\
\midrule
\endfirsthead
\toprule
\rowcolor{tableHeader}
\textbf{Kod} & \textbf{Historyjka} & \textbf{SP} & \textbf{Status} \\
\midrule
\endhead
S7-01 & Pełny przepływ klienta w UI & 8 & Done \\
S7-02 & Responsywność i spójność wizualna & 5 & Done \\
S7-03 & Rozszerzenie testów automatycznych backendu & 13 & Done \\
S7-04 & Kolekcja Bruno i testy manualne API & 5 & Done \\
S7-05 & Build, dokumentacja i scenariusz prezentacji & 5 & Done \\
\bottomrule
\end{longtable}
\endgroup

\subsection*{Podsumowanie sprintu}
Cel sprintu: Stabilizacja, testy i oddanie produktu. Dopracowanie całego przepływu klienta, usunięcie niespójności i przygotowanie produktu do prezentacji.

Zakres sprintu obejmował 5 historyjek o łącznej estymacji 36 SP. Wszystkie historyjki zostały dowiezione ze statusem Done, dlatego osiągnięta prędkość sprintu wyniosła 36 SP.

Przygotowano stabilną wersję produktu do prezentacji: pełny przepływ klienta w UI, rozszerzone testy backendu, uporządkowaną kolekcję Bruno, build frontendu oraz scenariusz demonstracyjny.

Sprint zamknął prace nad produktem w ramach zaplanowanego backlogu. Zespół skupił się na stabilności, jakości i prezentowalności rozwiązania, dzięki czemu aplikacja może zostać pokazana jako spójny przepływ bankowości internetowej.

\subsubsection*{Najważniejsze rezultaty}
\begin{itemize}
\item Zweryfikowano pełny przepływ klienta w UI.
\item Poprawiono responsywność i spójność wizualną.
\item Rozszerzono testy automatyczne backendu.
\item Uporządkowano kolekcję Bruno i scenariusz prezentacji.
\end{itemize}

\subsection*{Retrospekcja sprintu}
Retrospekcja końcowa pokazała, że regularne domykanie testów i dokumentacji w trakcie sprintów znacząco zmniejszyło koszt końcowej stabilizacji.

\subsubsection*{Co się udało}
\begin{itemize}
\item Udało się przejść pełny scenariusz klienta od rejestracji do dashboardu i przelewów.
\item Zespół uporządkował testy automatyczne i manualne.
\item Dokumentacja oraz Bruno ułatwiły przygotowanie prezentacji.
\item Poprawiono spójność UI i gotowość techniczną aplikacji.
\end{itemize}
\subsubsection*{Poziom energii zespołu}
Średni poziom energii zespołu: 4,2/5.

Energia zespołu wzrosła dzięki finalizacji produktu i widocznemu efektowi kilku sprintów pracy.

\subsubsection*{Czego udało się nauczyć}
\begin{itemize}
\item Zespół nauczył się stabilizować produkt przez połączenie testów, buildów i testów manualnych.
\item Scenariusz demo pomaga wykryć problemy integracyjne.
\item Regularne utrzymywanie dokumentacji zmniejsza koszt końcowego odbioru projektu.
\end{itemize}
\subsubsection*{Wnioski i działania na kolejny sprint}
\begin{itemize}
\item Zachować scenariusz demo jako checklistę regresji.
\item Przy dalszym rozwoju planować testy manualne i automatyczne razem z historyjkami.
\item Po odbiorze zebrać pomysły na kolejne epiki, np. historię operacji z paginacją, karty i przelewy cykliczne.
\end{itemize}

\subsection*{Następny sprint}
Product Owner / Scrum Master w następnym sprincie: Nie dotyczy - sprint 7 zamyka zakres produktu

Główny cel następnego sprintu: Nie zaplanowano kolejnego sprintu w ramach tego backlogu. Następnym krokiem jest prezentacja, odbiór produktu i ewentualne planowanie dalszego rozwoju.

\end{document}
